\chapter{Stratégie de données et coûts}\label{ch:fm:data-strategy-and-costs}

En février 2024, la Financial Conduct Authority britannique a publié les conclusions de son étude du marché de gros des données. Elle n'a trouvé aucune preuve que les
firmes ne pouvaient pas obtenir les données dont elles avaient besoin, et des signes de pouvoir de marché sur les trois marchés étudiés~: les indices de référence, les données de notation de crédit et les
fournisseurs de données de marché qui redistribuent les données de trading. Les utilisateurs, a-t-elle conclu, paient peut-être plus cher que ne le permettrait la concurrence~; la licence était
complexe, beaucoup d'utilisateurs devaient détenir des licences à la fois auprès d'une plateforme de négociation et auprès du fournisseur qui livrait leurs données, et cette complexité coûtait
elle-même de l'argent, jusque dans les équipes de conformité que les firmes emploient pour suivre leurs licences. La facture de données d'une société de trading dépend moins de la
quantité de données qu'elle utilise que de la façon dont elle les comptabilise.

\section{Ce qu'une firme paie pour les données}

\begin{definition}[Budget de données]\label{def:fm:data-strategy-and-costs:budget}
Un \emph{budget de données}\index{budget de données} est la dépense annuelle prévue d'une firme en données, par produit et par usage~: données de marché des bourses (affichage,
hors affichage et flux), terminaux et plateformes de fournisseurs, données de référence et d'opérations sur titres, actualités et jeux de données alternatives, avec les licences,
comptages et allocations qui sous-tendent chaque ligne.
\end{definition}

La firme de ce chapitre (\cref{tab:fm:data-strategy-and-costs:inventory}) dépense 2.07~millions de dollars par an en données~: 0.56~million de dollars en données de bourse, 0.75
millions de dollars pour trente terminaux de fournisseurs, 0.46~million de dollars en données de référence et actualités, et 0.30~million de dollars pour un jeu de données alternatives. Les terminaux
forment la plus grosse ligne, 36~\%, du total, et celle qui dépend le plus directement des effectifs. Les niveaux de redevance sont des données illustratives~; leur structure est ce que partagent
tous les contrats de données~: quelque chose de facturé par utilisateur, par appareil, par usage, ou un forfait, avec un plafond ou une alternative d'entreprise.

\begin{table}[ht]
\centering\small\setlength{\tabcolsep}{4pt}
\begin{tabular}{llrr}
\toprule
produit & facturation & usage & milliers de dollars par an \\
\midrule
affichage temps réel de bourse & 0.6 par utilisateur & 40 utilisateurs & 24 \\
bourse hors affichage (automatisé) & forfait & -- & 180 \\
flux directs & forfait, six flux & -- & 360 \\
terminaux de fournisseurs & 25 par utilisateur & 30 utilisateurs & 750 \\
données de référence & 8 par utilisateur ou 400 & 60 utilisateurs & 400 \\
analyse d'actualités & 15 par appareil & 4 appareils & 60 \\
jeu de données alternatives & forfait & -- & 300 \\
\midrule
total & & & 2\,074 \\
\bottomrule
\end{tabular}
\caption{L'inventaire des données du chapitre, chaque produit à sa licence la moins chère (niveaux de redevance illustratifs). Données~: \texttt{fm\_data}.}\label{tab:fm:data-strategy-and-costs:inventory}
\end{table}

\begin{omfigure}
\begin{tikzpicture}
\begin{axis}[width=10.5cm,height=4.8cm,xbar,bar width=10pt,xmin=0,xmax=900,xlabel={\small milliers de dollars par an},ytick={0,1,2,3},
  yticklabels={données de marché,terminaux,données de référence,données alternatives},y dir=reverse,ymin=-0.6,ymax=3.6,
  tick label style={font=\footnotesize},grid=major,grid style={gray!20},table/col sep=comma,nodes near coords,
  nodes near coords style={font=\scriptsize}]
\addplot[fill=omDef!55,draw=omDef] table[y=pos,x=cost]{figdata/desk/22-data-strategy-and-costs/budget.csv};
\end{axis}
\end{tikzpicture}

\omcaption{Le \omterm{def:fm:data-strategy-and-costs:budget}{budget de données} par catégorie, 2.07~millions de dollars par an~: les terminaux, ligne dictée par les effectifs, sont la plus grosse. Données~:
\texttt{fm\_data.the\_budget}.}\label{fig:fm:data-strategy-and-costs:budget}
\end{omfigure}

\section{Licences et leurs unités de comptage}

Le livre~15 a défini l'unité de comptage, le système de droits d'accès et le rapport d'utilisation~: la mécanique par laquelle une firme sait qui et quoi consomme chaque
jeu de données. Les licences qu'elles servent prennent quelques formes.

\begin{definition}[Licences d'entreprise, de données dérivées et de redistribution]\label{def:fm:data-strategy-and-costs:licences}
Une \emph{licence entreprise}\index{licence entreprise} permet un usage illimité d'un produit au sein de la firme (ou d'une partie définie de celle-ci) pour un montant forfaitaire fixe,
en remplacement des comptages par utilisateur ou par appareil. Une \emph{licence de données dérivées}\index{licence de données dérivées} permet à la firme de créer et d'utiliser, et
parfois de distribuer, des données calculées à partir des données sous licence, comme des prix, des indices ou des signaux. Une \emph{licence de redistribution
licence}\index{licence de redistribution} permet à la firme de transmettre les données, ou les données qui en dérivent, à des tiers tels que des clients, à certaines
conditions et moyennant redevances.
\end{definition}

La décision entre \omterm{def:fm:data-strategy-and-costs:licences}{licence entreprise} et licences par utilisateur est un seuil de rentabilité. Les données de référence coûtent 8\,000~\$ par utilisateur et par an, ou
400\,000~\$ pour l'entreprise~: la \omterm{def:fm:data-strategy-and-costs:licences}{licence entreprise} l'emporte à partir de 50 utilisateurs (\cref{fig:fm:data-strategy-and-costs:enterprise}). Avec les soixante utilisateurs de la firme, elle économise 80\,000~\$ par an, et elle supprime le
comptage et, avec lui, l'exposition à l'audit, pour ce produit. En dessous de cinquante utilisateurs, les redevances par utilisateur sont moins chères, sauf si les effectifs doivent dépasser le seuil
de rentabilité pendant la durée du contrat.

\begin{omfigure}
\begin{tikzpicture}
\begin{axis}[width=10.5cm,height=5cm,xlabel={\small utilisateurs},ylabel={\small milliers de dollars par an},xmin=0,xmax=100,ymin=0,ymax=820,
  tick label style={font=\footnotesize},grid=major,grid style={gray!20},table/col sep=comma,
  legend cell align=left,legend style={font=\scriptsize,at={(0.5,-0.3)},anchor=north,/tikz/every even column/.append style={column sep=8pt}},legend columns=2]
\addplot[omDef,thick] table[x=users,y=per_user]{figdata/desk/22-data-strategy-and-costs/enterprise.csv};
\addplot[omProp,thick,dashed] table[x=users,y=enterprise]{figdata/desk/22-data-strategy-and-costs/enterprise.csv};
\draw[gray,dotted] (axis cs:50,0) -- (axis cs:50,820);
\draw[gray,dotted] (axis cs:60,0) -- (axis cs:60,820);
\legend{redevances par utilisateur,licence entreprise}
\end{axis}
\end{tikzpicture}

\omcaption{Redevances par utilisateur face à une \omterm{def:fm:data-strategy-and-costs:licences}{licence entreprise} pour les données de référence~: les droites se croisent à 50 utilisateurs~; la firme en compte 60. Données~:
\texttt{fm\_data.reference\_curve}.}\label{fig:fm:data-strategy-and-costs:enterprise}
\end{omfigure}

\omcode{code/firm/databudget/firm_databudget.py}{36}{67}{Coût annuel d'un produit selon ses règles de redevance, budget à la licence la moins chère, seuil
de rentabilité de l'entreprise et exposition à un audit.}\label{lst:fm:data-strategy-and-costs:budget}

\section{Audits et facturation rétroactive}

\begin{definition}[Facturation rétroactive]\label{def:fm:data-strategy-and-costs:backbill}
La \emph{facturation rétroactive}\index{facturation rétroactive} est la somme qu'un fournisseur de données réclame après un audit pour des usages qui n'ont pas été déclarés ou payés
lors de périodes passées, généralement avec intérêts et parfois avec une pénalité, en vertu de la clause d'audit de sa licence.
\end{definition}

\begin{proposition}[Ce qu'un audit révèle]\label{prop:fm:data-strategy-and-costs:audit}
Si une firme a déclaré des redevances $F$ par an alors qu'une part $s$ de son usage réel est restée non déclarée pendant $M$ mois, les redevances non facturées valent $u=Fs/(12(1-s))$ par
mois, et la \omterm{def:fm:data-strategy-and-costs:backbill}{facturation rétroactive}, avec un intérêt mensuel $i$ depuis chaque mois jusqu'à l'audit, est
\[
u\sum_{k=0}^{M-1}(1+i)^k=u\,\frac{(1+i)^M-1}{i}.
\]
\end{proposition}

\begin{proof}
L'usage réel est l'usage déclaré divisé par $1-s$, donc les redevances non déclarées valent $F\,s/(1-s)$ par an. Le montant non facturé du mois $m$ porte intérêt pendant
$M-m$ mois~; la somme de la série géométrique donne la formule.
\end{proof}

Les frais par utilisateur de la société, affichage et terminaux compris, s'élèvent à 774\,000~\$ par an. Un audit qui constate que 15~\% des utilisateurs réels n'ont pas été déclarés sur trois ans facture
410\,000~\$ de frais et, à 1~\% par mois, 81\,000~\$ d'intérêts~: 490\,000~\$, soit 63~\% d'une année de frais par utilisateur, exigibles d'un seul coup
(la \cref{fig:fm:data-strategy-and-costs:audit}). L'exposition croît plus vite que la part~: une sous-déclaration de 5~\% coûte 146\,000~\$, une de 30~\% coûte
1.19~million de dollars. La parade est le système de droits d'accès du livre~15~: des comptages tirés des systèmes qui accordent l'accès, et non de ce que déclarent les utilisateurs.

\begin{omfigure}
\begin{tikzpicture}
\begin{axis}[width=10.5cm,height=5cm,ybar stacked,bar width=16pt,xlabel={\small part des utilisateurs réels non déclarés (\%)},ylabel={\small milliers de dollars},
  ymin=0,ymax=1300,xtick={5,10,15,20,25,30},tick label style={font=\footnotesize},grid=major,grid style={gray!20},table/col sep=comma,
  legend cell align=left,legend style={font=\scriptsize,at={(0.5,-0.3)},anchor=north,/tikz/every even column/.append style={column sep=8pt}},legend columns=2]
\addplot[fill=omDef!55,draw=omDef] table[x=under,y=principal]{figdata/desk/22-data-strategy-and-costs/audit.csv};
\addplot[fill=omProp!55,draw=omProp] table[x=under,y=interest]{figdata/desk/22-data-strategy-and-costs/audit.csv};
\legend{frais non facturés,intérêts à 1~\% par mois}
\end{axis}
\end{tikzpicture}

\omcaption{Rétrofacturation après un audit de trois années de frais par utilisateur (774\,000~\$ par an déclarés), selon la part des utilisateurs réels
qui n'ont pas été déclarés. Données~: \texttt{fm\_data.audit}.}\label{fig:fm:data-strategy-and-costs:audit}
\end{omfigure}

\section{Tutoriel~: la lettre d'audit}\label{tut:fm:data-strategy-and-costs:tut}

\begin{tutorial}
\textbf{Objectif.} Chiffrer l'inventaire de données d'une société, repérer où une \omterm{def:fm:data-strategy-and-costs:licences}{licence entreprise} devient avantageuse et estimer l'exposition d'un audit qui constate des utilisateurs
non déclarés. \textbf{État final~:} le budget par catégorie (la \cref{fig:fm:data-strategy-and-costs:budget}), le seuil de rentabilité de la \omterm{def:fm:data-strategy-and-costs:licences}{licence entreprise}
(la \cref{fig:fm:data-strategy-and-costs:enterprise}) et l'exposition à l'audit (la \cref{fig:fm:data-strategy-and-costs:audit}).
\begin{tutsteps}
\item \textbf{L'inventaire.} \texttt{fm\_data.PRODUCTS} et \texttt{USAGE} contiennent sept produits avec leurs règles de frais et leurs comptages~;
\texttt{firm.databudget.budget} chiffre chacun à sa licence la moins chère.
\item \textbf{Le seuil de rentabilité.} \texttt{enterprise\_breakeven} donne l'effectif à partir duquel la \omterm{def:fm:data-strategy-and-costs:licences}{licence entreprise} est moins chère~;
\texttt{fm\_data.reference\_curve} trace les deux coûts en fonction du nombre d'utilisateurs.
\item \textbf{L'audit.} \texttt{fm\_data.audit(under, years, monthly\_rate)} applique la \cref{prop:fm:data-strategy-and-costs:audit} aux
frais par utilisateur.
\item \textbf{La répartition.} \texttt{fm\_data.allocation()} répartit le budget entre quatre stratégies selon leur usage mesuré~;
\texttt{fm\_data.alt\_value()} chiffre le jeu de données alternatif avec l'arithmétique d'essai fournisseur du livre~7.
\end{tutsteps}
\end{tutorial}

\textbf{Pour aller plus loin.} Faire croître l'effectif de 20~\% par an et trouver quand les terminaux justifieraient une négociation entreprise~;
ajouter des frais de redistribution pour une page de cotation destinée aux clients.

\section{Négocier avec les bourses et les fournisseurs}

\begin{dated}{2026-09}{Règles sur la tarification des données de marché}\label{dat:fm:data-strategy-and-costs:rules}
\textbf{Union européenne}~: le MiFIR (règlement (UE) n\up{o}~600/2014), article~13, tel qu'adopté, exige que les plateformes de négociation rendent disponibles les données pré- et post-négociation
sur une base commerciale raisonnable, avec un accès non discriminatoire, et gratuitement quinze minutes après publication. \textbf{États-Unis}~:
les règles d'infrastructure des données de marché de la SEC prévoient des consolidateurs concurrents de données de marché NMS, enregistrés au titre de la Rule 614 (17 CFR
242.614, adoptée en 2021). \textbf{Royaume-Uni}~: l'étude de marché de la FCA sur les données de gros (février 2024) a relevé des indices de pouvoir de marché chez les fournisseurs d'indices de référence, les agences de notation de crédit et les fournisseurs de données de marché,
et a indiqué qu'elle poursuivrait ses constats par une réforme réglementaire et par ses pouvoirs
en matière de concurrence.
\end{dated}

La négociation part de l'inventaire et des comptages. Une société qui sait exactement qui utilise quoi peut faire passer des catégories entières à des conditions entreprise, retirer
les droits d'accès inutilisés avant le renouvellement et résister aux hausses de prix en présentant un plan crédible de changement de fournisseur là où un substitut existe. Elle dispose de moins de levier lorsque
les données proviennent d'une seule plateforme ou d'un seul administrateur, c'est-à-dire du pouvoir de marché décrit par la FCA~; là, les leviers portent sur les termes de la licence plutôt que sur le prix~:
définitions de l'usage non destiné à l'affichage, droits sur les données dérivées, périodes d'audit et intérêts, et plafonds sur les hausses au renouvellement.

\section{Les données comme actif stratégique}

Un jeu de données alternatif s'achète pour le rendement qu'il ajoute. Le harnais d'essai fournisseur du livre~7 le chiffre~: le montant annuel des frais auquel la contribution attendue
du jeu de données, qui s'érode à mesure que d'autres l'achètent, couvre tout juste son coût. Pour le jeu de données du chapitre, alimentant une stratégie sur 50~millions de dollars qui en tire 2~\% par an
avant des coûts de transaction de 0.5~\%, avec un edge qui se réduit de moitié tous les deux ans, les frais de seuil sur un contrat de trois ans s'élèvent à 369\,000~\$ par an
(\texttt{fm\_data.alt\_value})~: à 300\,000~\$, le jeu de données mérite d'être acheté, avec peu de marge d'erreur sur l'érosion.

Les coûts de données sont affectés aux stratégies selon leur usage~: avec les pondérations du chapitre, le market making porte 933\,000~\$ du budget, l'arbitrage statistique
622\,000~\$, le trading événementiel 311\,000~\$ et la macro 207\,000~\$. Une stratégie dont les coûts de données dépassent ce qu'elle en tire ne devrait pas
les porter, ce qu'on ne voit qu'une fois qu'ils sont affectés.

\begin{method}[Conduire un budget de données]\label{met:fm:data-strategy-and-costs:run}
\begin{enumerate}
\item Tenir un inventaire de chaque produit de données~: fournisseur, licence, unité de comptage, comptages actuels, date de renouvellement, responsable.
\item Tirer les comptages du système de droits d'accès, les rapprocher chaque mois des factures des fournisseurs et supprimer les accès inutilisés avant chaque renouvellement.
\item Comparer les conditions à l'unité et les conditions entreprise à l'effectif prévu~; chiffrer l'exposition à l'audit de chaque licence comptée.
\item Chiffrer les données alternatives par leur contribution attendue avant et après un essai, et arrêter ce qui ne rapporte pas.
\item Répartir le budget entre les stratégies selon leur usage et le revoir à la lumière de leurs résultats.
\end{enumerate}
\end{method}

\section{Construction~: le budget de données}\label{bld:fm:data-strategy-and-costs:databudget}

\begin{build}
\bfield{Rôle} Produits de données et règles de frais sous forme de données, le budget, le seuil de rentabilité entreprise, l'exposition à l'audit, la répartition entre les stratégies et la valeur
d'un jeu de données alternatif.
\bfield{Interface} \texttt{firm.databudget}~: \texttt{Product}, \texttt{Usage}, \texttt{annual\_cost}, \texttt{budget}, \texttt{enterprise\_breakeven},
\texttt{audit\_exposure}, \texttt{allocate}, \texttt{alt\_breakeven} (qui enveloppe \texttt{firm.vendoreval.breakeven}).
\bfield{Règles} Les niveaux de frais sont des entrées~; chaque produit est budgété à sa licence la moins chère~; la rétrofacturation accumule des intérêts de chaque mois jusqu'à l'audit.
\bfield{Tests d'acceptation} \texttt{code/firm/databudget/tests/}~: les coûts et le budget contre des chiffres calculés à la main~; le seuil de rentabilité~; la somme géométrique de l'audit~;
l'enveloppe.
\bfield{Pour aller plus loin} Enregistrements d'usage du système de droits d'accès du livre~15~; dates de renouvellement et préavis~; frais de données dérivées et de redistribution.
\end{build}

\begin{omsources}
\item Financial Conduct Authority, \emph{Wholesale Data Market Study}, Market Study MS23/1.5, février 2024.
\item Regulation (EU) No 600/2014, article 13~; 17 CFR 242.614.
\end{omsources}

\section{Exercices}

\begin{exercise}[$\star$]\label{exo:fm:data-strategy-and-costs:1}
À partir de quel effectif la \omterm{def:fm:data-strategy-and-costs:licences}{licence entreprise} des données de référence devient-elle avantageuse, et que fait-elle économiser à soixante utilisateurs~?
\end{exercise}

\begin{exercise}[$\star$]\label{exo:fm:data-strategy-and-costs:2}
Qu'a constaté l'étude de la FCA sur l'accès aux données et sur les prix~?
\end{exercise}

\begin{exercise}[$\star$]\label{exo:fm:data-strategy-and-costs:3}
Sans intérêts, combien un audit facture-t-il pour 15~\% des utilisateurs réels non déclarés sur trois ans, sur 774\,000~\$ par an de frais par utilisateur~?
\end{exercise}

\begin{exercise}[$\star\star$]\label{exo:fm:data-strategy-and-costs:4}
Démontrer la forme fermée de la \cref{prop:fm:data-strategy-and-costs:audit} et calculer les intérêts à 1~\% par mois.
\end{exercise}

\begin{exercise}[$\star\star$]\label{exo:fm:data-strategy-and-costs:5}
Expliquer la différence entre une \omterm{def:fm:data-strategy-and-costs:licences}{licence de données dérivées} et une \omterm{def:fm:data-strategy-and-costs:licences}{licence de redistribution} à partir d'un exemple tiré d'une société de trading.
\end{exercise}

\begin{exercise}[$\star\star$]\label{exo:fm:data-strategy-and-costs:6}
Pourquoi l'exposition à l'audit croît-elle plus vite que la part non déclarée~?
\end{exercise}

\begin{exercise}[$\star\star\star$]\label{exo:fm:data-strategy-and-costs:7}
\emph{Programmation.} Avec une demi-vie d'un an au lieu de deux, quels sont les frais d'équilibre du jeu de données alternatif~? Faut-il acheter à 300\,000~\$~?
\end{exercise}

\begin{exercise}[$\star\star\star$]\label{exo:fm:data-strategy-and-costs:8}
\emph{Trouver la faille.} «~Nos utilisateurs nous indiquent chaque trimestre quels terminaux ils utilisent, donc nos comptages sont justes.~»
\end{exercise}

\section{Problème~: la lettre d'audit}

\begin{problem}[{Problème du week-end --- la lettre d'audit}]\label{pb:fm:data-strategy-and-costs:1}
Un directeur des opérations reçoit l'avis d'audit d'une bourse et doit estimer l'exposition et repenser la manière dont la société achète des données.

\medskip\noindent\textbf{Partie I --- Le marché.}
\begin{enumerate}
\item Qu'a constaté la FCA en 2024, et pourquoi cela importe-t-il pour un acheteur~?
\item Qu'impose l'article 13 du MiFIR aux plateformes de négociation~?
\item Définir un \omterm{def:fm:data-strategy-and-costs:budget}{budget de données} et en énumérer les catégories.
\item Donner le budget de la société par catégorie.
\end{enumerate}

\medskip\noindent\textbf{Partie II --- Licences.}
\begin{enumerate}[resume]
\item Définir les \omterm{def:fm:data-strategy-and-costs:licences}{licences entreprise}, de données dérivées et de redistribution.
\item Calculer le seuil de rentabilité entreprise des données de référence.
\item Qu'est-ce qu'une \omterm{def:fm:data-strategy-and-costs:licences}{licence entreprise} change d'autre que le prix~?
\item Quelles clauses négocier lorsque le fournisseur détient un pouvoir de marché~?
\end{enumerate}

\medskip\noindent\textbf{Partie III --- L'audit.}
\begin{enumerate}[resume]
\item Définir la \omterm{def:fm:data-strategy-and-costs:backbill}{facturation rétroactive}.
\item Énoncer et démontrer la \cref{prop:fm:data-strategy-and-costs:audit}.
\item Donner l'exposition pour une sous-déclaration de 15~\% sur trois ans, puis pour 5~\% et 30~\%.
\item Comment répondre à l'avis~?
\item Comment un système de droits d'accès réduit-il l'exposition~?
\end{enumerate}

\medskip\noindent\textbf{Partie IV --- Stratégie.}
\begin{enumerate}[resume]
\item Comment tarifer le jeu de données alternatif, et quels sont ses frais d'équilibre~?
\item Répartir le budget entre les stratégies.
\item Qu'indique la répartition à la société~?
\item Quels achats arrêter, et comment trancher~?
\item À quelle fréquence faut-il réconcilier l'inventaire~?
\item Énoncer le \emph{résultat nommé}~: l'effectif à partir duquel une \omterm{def:fm:data-strategy-and-costs:licences}{licence entreprise} devient plus avantageuse que les frais par utilisateur, et l'exposition à la \omterm{def:fm:data-strategy-and-costs:backbill}{facturation rétroactive} d'un audit qui
constate une sous-déclaration de 15~\% sur trois ans.
\item En deux phrases, rédiger la politique de données.
\end{enumerate}
\end{problem}

\section{Questions d'entretien}

\begin{interviewq}[$\star$ \iqroles{développeur}]\label{iq:fm:data-strategy-and-costs:1}
Qu'est-ce qu'un usage non destiné à l'affichage des données de marché, et pourquoi importe-t-il pour une société de trading~?
\end{interviewq}

\begin{interviewq}[$\star$ \iqroles{trader}]\label{iq:fm:data-strategy-and-costs:2}
Pourquoi une société paierait-elle une \omterm{def:fm:data-strategy-and-costs:licences}{licence entreprise} qu'elle n'utilise pas entièrement~?
\end{interviewq}

\begin{interviewq}[$\star\star$ \iqroles{chercheur}]\label{iq:fm:data-strategy-and-costs:3}
Comment décideriez-vous si un jeu de données alternatif vaut son prix~?
\end{interviewq}

\begin{interviewq}[$\star\star$ \iqroles{développeur}]\label{iq:fm:data-strategy-and-costs:4}
Concevez un système qui garde les comptages de données de marché prêts pour un audit.
\end{interviewq}

\begin{interviewq}[$\star\star$ \iqroles{risque}]\label{iq:fm:data-strategy-and-costs:5}
Un audit de bourse constate une sous-déclaration. Quelle est l'exposition de la société, et que faites-vous ensuite~?
\end{interviewq}

\begin{interviewq}[$\star\star\star$ \iqroles{chercheur, trader}]\label{iq:fm:data-strategy-and-costs:6}
Comment les coûts de données doivent-ils être répartis entre les stratégies, et quel comportement la répartition engendre-t-elle~?
\end{interviewq}
